\section{Week 9: Evaluation Methods for AI Systems}

\subsection{Topic Summary}

\begin{keybox}[Learning Objectives]
By the end of this week, you will be able to:
\begin{itemize}
  \item Understand the range of methods available to a human-AI designer/practitioner, and describe their pros/cons.
  \item Apply both simple methods (think-aloud protocol, heuristic analysis) and experiment design thinking.
  \item Run a think-aloud protocol and extract useful feedback for designs.
  \item Reason about independent, dependent, and intervening variables and argue about causal effects of system design.
\end{itemize}
\end{keybox}

\begin{keybox}[Big Picture]
Evaluation methods are your \textbf{toolbox}---each has a purpose, and you must know when/why to use each. This week covers:
\begin{enumerate}
  \item \textbf{Heuristic evaluation}: Expert-based, uses established principles.
  \item \textbf{Think-aloud protocol}: Users verbalize thoughts while using system.
  \item \textbf{Experiments}: Manipulate variables to establish causation.
\end{enumerate}

Key insight: \textbf{Correlation $\neq$ Causation}. Experiments help us make causal claims.
\end{keybox}

\begin{keybox}[What Examiners Like to Ask]
\begin{itemize}
  \item Compare heuristic evaluation and think-aloud protocol.
  \item Explain independent vs.\ dependent vs.\ intervening variables.
  \item Identify threats to internal/external validity and reliability.
  \item Design an experiment with appropriate controls.
  \item Compare between-subjects vs.\ within-subjects designs.
  \item Explain why triangulation is important.
\end{itemize}
\end{keybox}

\subsection{Overview + Core Concepts + Extended Theory}

\subsubsection*{Evaluation in the Design Process}

\paragraph*{Why Evaluate?}
To check that all aspects of user experience are considered in design.

\paragraph*{What to Evaluate?}
\begin{itemize}
  \item Conceptual models
  \item Early and subsequent prototypes
  \item Complete prototypes
  \item Comparisons with competitors' products
\end{itemize}

\paragraph*{Where?}
Natural settings, in-the-wild, living labs, lab settings.

\paragraph*{When?}
Throughout design---even finished products can be evaluated.

\subsubsection*{Map of Evaluation Methods}

\begin{center}
\renewcommand{\arraystretch}{1.3}
\begin{tabular}{@{} l p{9cm} @{}}
\toprule
\textbf{Category} & \textbf{Methods} \\
\midrule
\textbf{Expert evaluations} (no users) & Heuristic evaluation, model-based evaluation, analytics, cognitive walkthrough \\
\textbf{Lab studies} & Think-aloud protocol, experiments, usability evaluation \\
\textbf{Surveys} & Questionnaires, structured interviews \\
\textbf{Field studies} & Technology probes, prototype deployment, A/B testing \\
\textbf{Living laboratories} & Long-term instrumented environments \\
\bottomrule
\end{tabular}
\end{center}

\paragraph*{Three Dimensions of Methods}
\begin{enumerate}
  \item \textbf{Abstract vs.\ Concrete}: Are you testing an abstraction or actual user behavior?
  \item \textbf{Obtrusive vs.\ Unobtrusive}: Do users come to you (lab) or do you observe naturally?
  \item \textbf{Goals}: Realistic use? Generalizability? Precision?
\end{enumerate}

\subsubsection*{Method 1: Heuristic Evaluation}

\begin{defbox}[Heuristic Evaluation]
Experts evaluate whether UI elements conform to tried-and-tested usability principles (heuristics). It is conducted \textbf{by experts}, not users.
\end{defbox}

\paragraph*{Nielsen's 10 Heuristics (Review)}
\begin{enumerate}
  \item Visibility \& Feedback
  \item Match between system \& real world
  \item User control \& freedom
  \item Consistency \& standards
  \item Error prevention
  \item Recognition rather than recall
  \item Efficiency \& flexibility
  \item Aesthetics \& minimalism
  \item Help users recognize, diagnose, recover from errors
  \item Help and documentation
\end{enumerate}

\paragraph*{Process}
\begin{enumerate}
  \item \textbf{Recruit experts}: People who understand heuristics, have design experience.
  \item \textbf{Briefing}: Prepared script ensures everyone receives same briefing; provide specific user tasks to focus exploration.
  \item \textbf{Evaluation}: At least \textbf{two passes}:
  \begin{itemize}
    \item Pass 1: Get feel for flow and scope.
    \item Pass 2: Focus on specific elements, identify usability problems.
  \end{itemize}
  \item \textbf{Debrief}: Summarize issues and recommendations to designers.
\end{enumerate}

\begin{tipbox}[When to Use Heuristic Evaluation]
\begin{itemize}
  \item Early in design (on sketches, wireframes).
  \item When experts are available but users are not.
  \item As a ``cheap, quick'' first pass before user testing.
  \item Surface-level check regardless of backend functionality.
\end{itemize}
\end{tipbox}

\subsubsection*{Method 2: Think-Aloud Protocol}

\begin{defbox}[Think-Aloud Protocol]
Users use the system while \textbf{continuously verbalizing their thoughts}---what they're thinking, trying to do, questions that arise, things they read.
\end{defbox}

\paragraph*{Why?}
Need to know what users are \textbf{thinking}, not just what they're \textbf{doing}.

\paragraph*{Process}
\begin{enumerate}
  \item Give user specific tasks (core features you want to test).
  \item Prompt: ``Tell me what you're thinking.''
  \item \textbf{Only help on pre-decided things}---track anything you do help on.
  \item Make recording \& take good notes:
  \begin{itemize}
    \item Screen capture
    \item Audio/video
    \item Timestamps
    \item Event logs
  \end{itemize}
  \item Do \textbf{not} answer questions or guide user (except pre-decided help).
\end{enumerate}

\paragraph*{What to Capture}
\begin{itemize}
  \item What are they thinking?
  \item What are they trying to do?
  \item What questions arise?
  \item What are they reading/looking at?
  \item Where do they pause/hesitate?
\end{itemize}

\begin{tipbox}[Think-Aloud Best Practices]
\begin{itemize}
  \item Be quiet---only prompt ``tell me what you're thinking.''
  \item Don't answer their questions (you're testing the system, not your guidance).
  \item Record everything for later analysis.
  \item Note where users are confused vs.\ just thinking hard.
\end{itemize}
\end{tipbox}

\subsubsection*{Method 3: Experiments}

\paragraph*{The Problem: Correlation $\neq$ Causation}

\begin{exbox}[Classic Example]
Ice cream sales and shark attacks both increase in summer. Does ice cream attract sharks? No---both are caused by a \textbf{third variable} (more people at the beach in warm weather).
\end{exbox}

Experiments help us separate \textbf{accidental correlations} from \textbf{causal relationships}.

\begin{defbox}[Experiment]
A controlled study that:
\begin{enumerate}
  \item States a \textbf{testable hypothesis} (if X then Y).
  \item \textbf{Manipulates independent variables}.
  \item \textbf{Measures dependent variables}.
  \item Uses statistical tests to accept/reject hypothesis.
\end{enumerate}
\end{defbox}

\paragraph*{Key Variables}

\begin{center}
\renewcommand{\arraystretch}{1.3}
\begin{tabular}{@{} l p{7cm} l @{}}
\toprule
\textbf{Variable Type} & \textbf{Definition} & \textbf{Example} \\
\midrule
Independent (X) & What you \textbf{manipulate} & Mac vs.\ Windows menu bar \\
Dependent (Y) & What you \textbf{measure} & Time to access menu, errors \\
Intervening ($\varepsilon$) & Unknown/uncontrolled factors & User fatigue, prior experience \\
\bottomrule
\end{tabular}
\end{center}

\textbf{Model}: $Y = f(X) + \varepsilon$

The purpose of experiment design is to \textbf{minimize or control} $\varepsilon$ so we can draw conclusions about how X affects Y.

\begin{exbox}[Menu Bar Experiment]
\textbf{Hypothesis}: Mac menu bar (anchored to top of screen) is faster to access than Windows menu bar (attached to application window).

\textbf{Independent variables}:
\begin{itemize}
  \item Interface style (Mac vs.\ Windows)
  \item Training given to users
\end{itemize}

\textbf{Dependent variables}:
\begin{itemize}
  \item Time to access menu
  \item Errors
  \item Number of tasks completed
  \item User satisfaction
\end{itemize}

\textbf{Intervening variables} (must control):
\begin{itemize}
  \item Pointing device (mouse vs.\ trackpad)
  \item User fatigue
  \item Prior familiarity with each OS
\end{itemize}
\end{exbox}

\subsubsection*{Threats to Validity and Reliability}

\begin{center}
\renewcommand{\arraystretch}{1.3}
\begin{tabular}{@{} l p{5cm} p{5cm} @{}}
\toprule
\textbf{Threat} & \textbf{Question} & \textbf{Risk} \\
\midrule
Internal validity & Are results caused by independent variables? & Confounding variables cause effect, not your manipulation \\
External validity & Can results be generalized? & Lab conditions too artificial to apply to real world \\
Reliability & Will results be consistent if repeated? & Random variation, single trials, non-replicable \\
\bottomrule
\end{tabular}
\end{center}

\paragraph*{Threats to Internal Validity}

\begin{center}
\renewcommand{\arraystretch}{1.3}
\begin{tabular}{@{} l p{5cm} p{5cm} @{}}
\toprule
\textbf{Threat} & \textbf{Problem} & \textbf{Solution} \\
\midrule
Ordering effects & People learn, get tired & \textbf{Randomize} or \textbf{counterbalance} order \\
Selection effects & Pre-existing differences between groups & \textbf{Randomly assign} users to conditions \\
Experimenter bias & Enthusiasm for one condition & \textbf{Equivalent training}; \textbf{double-blind} design \\
\bottomrule
\end{tabular}
\end{center}

\paragraph*{Threats to External Validity}

\begin{center}
\renewcommand{\arraystretch}{1.3}
\begin{tabular}{@{} l p{5cm} p{5cm} @{}}
\toprule
\textbf{Threat} & \textbf{Problem} & \textbf{Solution} \\
\midrule
Population & Sample not representative & Draw \textbf{random sample} from target population \\
Ecological & Lab conditions unrealistic & Make conditions \textbf{realistic in important respects} \\
Training & Training doesn't match real learning & Mimic \textbf{real interface encounters} \\
Task & Tasks not representative & Base on \textbf{task analysis} \\
\bottomrule
\end{tabular}
\end{center}

\subsubsection*{Experiment Design: Between vs.\ Within Subjects}

\begin{center}
\renewcommand{\arraystretch}{1.3}
\begin{tabular}{@{} l p{5.5cm} p{5.5cm} @{}}
\toprule
& \textbf{Between Subjects} & \textbf{Within Subjects} \\
\midrule
\textbf{Design} & Different groups see different conditions & Same participants see all conditions \\
\textbf{Comparison} & Compare means across groups & Compare within each participant \\
\textbf{Advantages} & No ordering/training effects & Fewer participants needed; controls for individual differences \\
\textbf{Disadvantages} & Needs 2$\times$+ participants; inter-participant variation can swamp effects & Risk of ordering/training effects; must counterbalance \\
\bottomrule
\end{tabular}
\end{center}

\begin{exbox}[Between Subjects]
\begin{itemize}
  \item Group A sees only Mac menu bar.
  \item Group B sees only Windows menu bar.
  \item Compare: Is mean(time for A) < mean(time for B)?
\end{itemize}
\textbf{Eliminates}: Learning from one interface to the other.

\textbf{Problem}: User differences (skill, preference) may swamp effect.
\end{exbox}

\begin{exbox}[Within Subjects]
\begin{itemize}
  \item Half see Mac first, then Windows.
  \item Half see Windows first, then Mac.
  \item Compare: For each person, was Mac faster than Windows?
\end{itemize}
\textbf{Eliminates}: Inter-participant variation (each person compared to self).

\textbf{Problem}: Learning/fatigue effects; must counterbalance order.
\end{exbox}

\subsubsection*{Counterbalancing}

\begin{defbox}[Counterbalancing]
Systematically varying the order of conditions across participants to control for ordering effects.
\end{defbox}

\textbf{Example} (2 conditions):
\begin{itemize}
  \item Half of participants: A $\to$ B
  \item Half of participants: B $\to$ A
\end{itemize}

\textbf{Note}: Becomes complex with $>2$ conditions (Latin square designs).

\subsubsection*{Triangulation}

\begin{defbox}[Triangulation]
Using \textbf{several different methods} to attack the same question and see if results concur or conflict.
\end{defbox}

\textbf{Why?} Every method has limitations:
\begin{itemize}
  \item Lab experiment: Not externally valid.
  \item Field study: Not controlled.
  \item Survey: Biased by self-report.
\end{itemize}

\textbf{Higher confidence}: Results from multiple methods pointing in same direction.

\subsubsection*{Comparison: Heuristic vs.\ Think-Aloud vs.\ Experiment}

\begin{center}
\renewcommand{\arraystretch}{1.3}
\begin{tabular}{@{} l c c c @{}}
\toprule
& \textbf{Heuristic} & \textbf{Think-Aloud} & \textbf{Experiment} \\
\midrule
Participants & Experts & Users & Users \\
Setting & Lab/desk & Lab & Lab \\
Goal & Check against heuristics & Understand user thinking & Establish causation \\
Output & Usability problems & Mental model insights & Statistical evidence \\
Cost & Low & Medium & High \\
When & Early prototypes & Interactive prototypes & Near-final designs \\
\bottomrule
\end{tabular}
\end{center}

\subsubsection*{Practical Application Flow}

A typical evaluation progression:
\begin{enumerate}
  \item \textbf{Heuristic evaluation} on semi-interactive sketches $\to$ check UI is understandable.
  \item \textbf{Think-aloud protocol} on clickable prototype $\to$ understand user mental models.
  \item \textbf{Small experiment} on near-final version $\to$ test if AI scaffolding works.
  \item \textbf{Technology probe/field study} $\to$ observe real-world use over days/weeks.
\end{enumerate}

\subsubsection*{Key Definitions Summary}

\begin{center}
\renewcommand{\arraystretch}{1.3}
\begin{tabular}{@{} l p{8cm} @{}}
\toprule
\textbf{Term} & \textbf{Definition} \\
\midrule
Heuristic evaluation & Expert review against usability principles \\
Think-aloud protocol & Users verbalize thoughts while using system \\
Independent variable & What experimenter manipulates \\
Dependent variable & What experimenter measures \\
Intervening variable & Uncontrolled factors affecting results \\
Internal validity & Results caused by independent variables \\
External validity & Results generalizable to real world \\
Reliability & Results repeatable \\
Confounding & Uncontrolled variable systematically affects results \\
Between subjects & Different groups see different conditions \\
Within subjects & Same participants see all conditions \\
Counterbalancing & Varying order of conditions across participants \\
Triangulation & Using multiple methods on same question \\
\bottomrule
\end{tabular}
\end{center}

\subsubsection*{Exam Tips}

\begin{tipbox}[Experiment Design Questions]
When asked to design an experiment:
\begin{enumerate}
  \item State the \textbf{hypothesis} clearly.
  \item Identify \textbf{independent variables} (what you manipulate).
  \item Identify \textbf{dependent variables} (what you measure).
  \item List \textbf{intervening variables} to control.
  \item Choose \textbf{between/within subjects} and justify.
  \item Explain how you'll address \textbf{validity threats}.
\end{enumerate}
\end{tipbox}


